% Page 050 — page imprimée 47
% Contenu seul : le préambule, l'en-tête et la pagination sont gérés par meyrueis.tex

\begin{center}
{\Large\bfseries Les chantiers de jeunesse à Meyrueis\par}
{\Large\bfseries 1940 -- 1944\par}
\end{center}

Les \textbf{Chantiers de la jeunesse française}, souvent appelés \textbf{chantiers de jeunesse}, étaient une organisation
paramilitaire française ayant existé de 1940 à 1944.

\textbf{\Large Histoire}

L'armistice de 1940 ayant supprimé le service militaire obligatoire, les chantiers de jeunesse furent créés
comme une sorte de substitut le 30 juillet 1940. Les jeunes hommes de la zone libre et de l'Afrique du nord
française en âge d'accomplir leur obligations militaires y étaient incorporés pour un stage de huit mois. Ils
vivaient en camps près de la nature, à la manière du scoutisme, et accomplissaient des travaux d'intérêt
général, notamment forestiers, dans une ambiance militaire. Ils étaient encadrés par des officiers d'active et de
réserve, ainsi que par des aspirants formés pendant la guerre de 1939-1940.

Dirigés par le général Paul de La Porte du Theil, les chantiers de jeunesse étaient une institution ambiguë. Il
s'agissait d'inculquer les valeurs de la Révolution nationale, prônée par le Régime de Vichy. Les Chantiers
étaient interdits aux Juifs aussi bien en métropole qu'en Afrique du Nord. Le culte de la hiérarchie et de la
discipline passait notamment par l'importance donnée à tous niveaux au Chef (toujours écrit avec une
majuscule). La vénération du Maréchal Pétain imprégnait profondément les cadres. Alors que le régime
exaltait le retour à la terre et le provincialisme, la vie en groupements dans les bois pouvait aussi se lire
comme une réaction contre la ville industrielle et corruptrice, foyer de l'individualisme et de la lutte des
classes.

Surtout après la guerre, les dirigeants des CJF affirmeront avoir voulu préparer une troupe mobilisable en cas
de reprise de la guerre contre l'Allemagne. Cette affirmation paraît discutable.

Le général La Porte du Theil choisit en effet de revenir d'Afrique du Nord en novembre 1942 au moment
même du débarquement anglo-saxon, au lieu de rejoindre le camp allié comme le firent les officiers vichystes
conduits par Darlan et Giraud. Et les Allemands se gardent de dissoudre les Chantiers lorsqu'ils envahissent la
Zone Sud, et qu'ils suppriment l'Armée d'Armistice. Significativement, l'occupant n'a toujours appelé les
Chantiers que "fransözische Arbeitsdienst", "service français du travail".

En 1943, les chantiers deviennent un piège pour des milliers de jeunes gens astreints au Service du Travail
Obligatoire (STO). Nombre d'entre eux sont envoyés en Allemagne par groupements entiers, avec le concours
actif de La Porte du Theil et de ses subordonnés. Ce concours est étayé tant par les faits et les témoignages
que par leurs instructions secrètes, conservées aux Archives Nationales [cote F60 1452 p.ex.]. On vit les
jeunes partir et arriver au travail forcé outre-Rhin en uniforme, en rangs, bannières déployées, encadrés par
leurs Chefs, et parfois musique en tête.

De nombreux jeunes s'enfuirent à temps des Chantiers. Plusieurs anciens des chantiers rejoignirent la
Résistance, certains passèrent en Afrique du Nord. Les stocks de vivres et de vêtements des Chantiers
devinrent une proie classique pour les maquis en manque de tout. De multiples coups de main, bénéficiant
parfois de complicités dans la place, permirent aux maquisards de faire main-basse sur ces dépôts - ce qui
explique que sur nombre de photos, des maquisards soient habillés en uniformes des chantiers, certes pas
prévus au départ pour leur servir !

Ce n'est qu'à partir de fin 1943 que le général La Porte du Theil, devant de nouvelles exigences allemandes
qui auraient signifié l'envoi en Allemagne de la quasi-totalité des effectifs restant des chantiers, commence
à se rebiffer. Il refuse toutefois les offres de la Résistance à passer à la dissidence et à gagner Alger. Averti
d'une arrestation imminente, il refuse également de s'enfuir. Arrêté le 4 janvier 1944 à son bureau de Châtel-
Guyon, il est assigné à résidence en Allemagne jusqu'à la fin de la guerre, épisode qui lui permit d'obtenir un
non-lieu en Haute-Cour.

\begin{flushright}
Extraits de l'encyclopédie Wikipédia
\end{flushright}
